\chapter{Contexto}
\label{ch:contexto}

\section{Inteligencia Artificial Generativa}

La Inteligencia Artificial (IA) Generativa agrupa a los modelos capaces de producir
contenido nuevo --texto, código, imágenes-- a partir de un patrón aprendido de
grandes cantidades de datos, en lugar de limitarse a clasificar o predecir sobre
datos ya existentes. Los modelos de lenguaje son el caso particular relevante para
este trabajo: reciben texto como entrada y generan texto como salida, condicionados
por una instrucción (\textit{prompt}).

\section{Modelos de Lenguaje de Gran Tamaño}

Los Modelos de Lenguaje de Gran Tamaño (LLM, \textit{Large Language Models}) son
redes neuronales entrenadas sobre corpus masivos de texto que, dada una secuencia de
entrada, predicen la continuación más probable. Esta capacidad, combinada con un
ajuste posterior para seguir instrucciones, permite emplearlos para tareas tan
diversas como resumir, traducir, generar código o --el uso que se les da en este
TFG-- extraer información estructurada de texto libre y redactar respuestas a
partir de datos ya recuperados.

Un LLM no tiene, por sí mismo, acceso a información posterior a su entrenamiento
ni a documentos privados como las actas de un ayuntamiento concreto. Tampoco
garantiza que lo que genera sea correcto: puede producir afirmaciones plausibles
pero falsas, fenómeno conocido como \textit{alucinación}. Ambas limitaciones son
las que motivan la siguiente técnica.

\section{Retrieval Augmented Generation (RAG)}

Retrieval Augmented Generation (RAG) es una arquitectura que combina un sistema de
recuperación de información con un LLM: antes de generar la respuesta, se buscan
en una base documental los fragmentos más relevantes para la pregunta y se
proporcionan al modelo como contexto. El LLM genera entonces la respuesta
apoyándose únicamente en esos fragmentos, no en su conocimiento interno.

Esto tiene dos ventajas directas para este proyecto: permite responder sobre un
corpus privado y actualizable (las actas del Pleno) sin reentrenar ningún modelo,
y permite citar la fuente exacta de cada afirmación, ya que se conoce qué
fragmentos --y por tanto qué documento y página-- sustentan la respuesta.

La recuperación de los fragmentos relevantes suele apoyarse en \textit{embeddings}:
representaciones vectoriales del texto en las que fragmentos semánticamente
similares quedan cerca en el espacio vectorial, lo que permite buscar por
significado y no solo por coincidencia literal de palabras.

\section{Grafos de Conocimiento}

Un grafo de conocimiento representa la información como una red de entidades
(nodos) y relaciones entre ellas (aristas), en lugar de como texto libre. El
estándar utilizado en este trabajo es RDF (\textit{Resource Description
Framework}), que modela cada hecho como una tripleta \textit{sujeto-predicado-objeto}
(por ejemplo, \texttt{proposición~X --presentadaPor--> Grupo~PP}). Sobre RDF se
puede definir una ontología con OWL (\textit{Web Ontology Language}) --las clases y
relaciones válidas del dominio-- y consultar el grafo con el lenguaje SPARQL.

Una ventaja de este modelo frente al texto libre es el razonamiento automático: un
motor de inferencia (\textit{reasoner}) puede derivar hechos nuevos que no están
escritos explícitamente pero se siguen de la ontología --por ejemplo, que una
proposición sobre "desahucios" trata también, por herencia temática, de
"vivienda"-- sin que haga falta repetir esa relación en cada dato.

\section{GraphRAG}

GraphRAG es la variante de RAG en la que la fuente de recuperación no es un
almacén de fragmentos de texto sino un grafo de conocimiento consultado mediante
un lenguaje de consulta estructurado (SPARQL en este trabajo). En lugar de
recuperar pasajes de texto parecidos a la pregunta, un LLM traduce la pregunta en
lenguaje natural a una consulta formal, el motor de grafos la ejecuta y devuelve
filas de datos exactos, que un segundo LLM convierte en una respuesta narrada.

Esto resulta especialmente adecuado para preguntas de agregación y razonamiento
multi-salto --"¿qué grupo ha presentado más proposiciones sobre vivienda?",
"¿cuántas se han aprobado por año?"-- que un sistema de recuperación de texto
libre no puede responder con precisión, porque exigen contar y cruzar datos de
todo el corpus, no solo encontrar un pasaje parecido a la pregunta. Es, junto con
el RAG vectorial tradicional, uno de los dos sistemas comparados en este TFG.
